\section{Luật chơi và Cơ chế tính điểm}

\subsection{Các loại khối gạch (Tetrominoes)}
Tetromino là hình khối được tạo thành từ 4 ô vuông ghép lại với nhau theo các cạnh.
Trong Tetris có đúng 7 loại tetromino chuẩn: \textbf{I, O, T, S, Z, J, L}.

\subsubsection{Hình dạng ma trận 4$\times$4 của từng khối}

\begin{center}
\begin{tabular}{ccc}
\textbf{I} & \textbf{O} & \textbf{T} \\
\texttt{....} & \texttt{.XX.} & \texttt{.X..} \\
\texttt{XXXX} & \texttt{.XX.} & \texttt{XXX.} \\
\texttt{....} & \texttt{....} & \texttt{....} \\
\texttt{....} & \texttt{....} & \texttt{....} \\
\end{tabular}
\end{center}

\begin{center}
\begin{tabular}{ccc}
\textbf{S} & \textbf{Z} & \textbf{J} \\
\texttt{.XX.} & \texttt{XX..} & \texttt{X...} \\
\texttt{XX..} & \texttt{.XX.} & \texttt{XXX.} \\
\texttt{....} & \texttt{....} & \texttt{....} \\
\texttt{....} & \texttt{....} & \texttt{....} \\
\end{tabular}
\end{center}

\begin{center}
\begin{tabular}{c}
\textbf{L} \\
\texttt{..X.} \\
\texttt{XXX.} \\
\texttt{....} \\
\texttt{....} \\
\end{tabular}
\end{center}

\subsubsection{Bảng màu sắc, ký tự đại diện và đặc điểm xoay của các Tetromino}

\begin{table}[H]
\centering
\begin{tabular}{|c|c|c|c|p{5cm}|}
\hline
\textbf{Khối} & \textbf{Ký tự} & \textbf{Màu chuẩn} & \textbf{Số trạng thái xoay} & \textbf{Ghi chú} \\
\hline
I & I & Cyan & 2 & Đối xứng qua tâm \\
\hline
O & O & Vàng & 1 & Không đổi hình khi xoay \\
\hline
T & T & Tím & 4 & Không đối xứng \\
\hline
S & S & Xanh lá & 2 & Đối xứng qua tâm \\
\hline
Z & Z & Đỏ & 2 & Đối xứng qua tâm \\
\hline
J & J & Xanh dương & 4 & Ảnh gương của khối L \\
\hline
L & L & Cam & 4 & Ảnh gương của khối J \\
\hline
\end{tabular}
\caption{Đặc điểm 7 loại Tetromino}
\end{table}

\subsection{Quy tắc xếp hình và Xóa dòng}
% [Gợi ý nội dung: Quy tắc rơi từ đỉnh xuống đáy, va chạm, ghép nối lấp đầy các hàng ngang để ghi điểm, kết thúc khi khối gạch chạm nóc (Game Over)]
\textit{[Nội dung bổ sung sau: Chi tiết quy tắc xếp và xử lý va chạm, kết thúc lượt chơi]}

\subsection{Hệ thống tính điểm và Cấp độ (Level / Score System)}
% [Gợi ý nội dung: Bảng điểm khi xóa 1 dòng (Single), 2 dòng (Double), 3 dòng (Triple), 4 dòng (Tetris), cơ chế tăng tốc độ rơi theo level]
\textit{[Nội dung bổ sung sau: Bảng tính điểm và công thức tăng độ khó]}
